\section{Problem Formulation: Automated Failure Attribution in Multi-Agent Cooperation}
\label{sec:problem_formulation}

In this section, we introduce decisive errors and formulate the automated failure attribution problem. We adopt the widely-adopted turn-based LLM multi-agent protocol~\citep{wu2023autogen,li2023camel,hong2023metagpt}.

\paragraph{Background.} Considering a LLMs-powered multi-agent system $\mathcal{M}$ with a group of $N$ agents, denoted as $\mathcal{N} = \{1, 2, .. N\}$, 
that operate at discrete time steps. These agents are taking actions in a turn-based protocol, meaning that exactly one agent performs an action at each time step. Formally, the system is described as:
\begin{align}
\mathcal{M} 
\;=\; 
\Bigl\langle 
  \mathcal{N},\, 
  S,\, 
  A,\, 
  P,\, 
  \phi 
\Bigr\rangle.
\end{align}
Here, $S$ is the set of possible states. $A$ is the global action set; each agent $i\in\mathcal{N}$ can typically perform actions from some subset $A_i \subseteq A$.
$\phi(t)$ is a function that indicates which agent is active at time $t$, thus specifying the turn-based rule. $P\bigl(s_{t+1} \mid s_t,\, a_t,\, \phi(t)\bigr)$ is the state-transition probability, given that \emph{only one} agent $\phi(t)$ acts at time $t$.

We employ $\phi(t)$ to denote the agent that takes an action $a_t$ at time step $t$.  A full trajectory $\tau$ can be written as:
$
\tau 
\;=\;
\bigl( s_0,\, a_0,\,
        s_1,\, a_1,\,
        \dots,\,
        s_T \bigr),
$
where $T$ is a terminal time step or when the system enters a terminating state.


\paragraph{Decisive Error and Objective.}
We use a tuple $(i,t)$ to denote a mistake in a trajectory, which means agent $i$ is active at time $t$, and its action $a_t$ is deemed a mistake (e.g., wrong reasoning etc.). 
A trajectory may contain multiple mistakes, but not all of them result in overall failure.
We employ $Z(\tau)$ to denote the result of a trajectory $\tau$. 
\begin{align}
Z(\tau)
\;=\;
\begin{cases}
1, 
& 
\text{if the system ultimately fails},\\[6pt]
0, 
& 
\text{otherwise}.
\end{cases}
\end{align}
Suppose the original trajectory $\tau$ is a failure, i.e., $Z(\tau)=1$.
Considering the following scenario, if correcting the mistake made by agent $i$ at time $t$: we replace $a_t$ with a "correct" action $\tilde{a}_t$. The steps prior to step $t$ remain unchanged, while the actions following $t$ are adjusted accordingly to ensure correctness. This process generates a modified trajectory:
\begin{align}
    \tau^{(i,t)} 
    \;=\;
    \mathcal{I}_{(i,t)}(\tau),
\end{align}
where $\mathcal{I}_{(i,t)}$ denotes the intervention.
If in the modified trajectory we obtain $Z\bigl(\tau^{(i,t)}\bigr)=0$ (success), then the error $(i,t)$ is said to be a \emph{decisive error}.
Formally, we define the \emph{decisive error indicator} $\Delta_{i,t}(\tau)$ as
\begin{align}
\Delta_{i,t}(\tau)
\;=\;
\begin{cases}
1, 
& 
\text{if } Z(\tau) = 1 \text{ and } Z\bigl(\tau^{(i,t)}\bigr) = 0,\\[6pt]
0, 
& 
\text{otherwise}.
\end{cases}
\end{align}
In words,
$
\Delta_{i,t}(\tau) = 1
\quad
\Longleftrightarrow
\quad$
Fixing agent $i$'s mistake at time $t$ changes $Z(\tau)$ from $1$ (fail) to $0$ (success). Formally, the decisive error could be defined as agent-time pairs $(i^*, t^*)$ such that $\Delta_{i^*,t^*}(\tau) = 1$, where $i^*$ represents the agent responsible for the system failure, and $t^*$ represents the exact time step at which the critical mistake occurs. We refer to these as the \emph{failure-responsible agent} and the \emph{decisive error step}, respectively across the paper.

In practice, multiple decisive errors may occur within a trajectory. In our study, we address this situation by identifying the earliest error in time as the principal cause of failure. Specifically, we define an objective to determine:
\begin{align}
\mathcal{C}(\tau) 
&=\; 
\bigl\{
  (i,t) \;\mid\; \Delta_{i,t}(\tau) \;=\; 1
\bigr\},
\label{eq:culprit_set}
\\[5pt] 
&(i^*,\, t^*)
=\;
\arg\min_{(i,t)\,\in\,\mathcal{C}(\tau)} 
\; t. \nonumber
\label{eq:earliest_culprit}
\end{align}
which selects the pair $(i^*, t^*)$ yielding the highest decisive error indicator with earliest time step.
In this study, the research problem focuses on the automatic identification of the $(i^*, t^*)$ in LLMs-powered multiple agent systems. 

